\documentclass[a4paper,12pt]{article}
\usepackage[utf8]{inputenc}
\usepackage[T1]{fontenc}
\usepackage[ngerman]{babel}
\usepackage{geometry}
\usepackage{amsmath}
\usepackage{hyperref}
\geometry{margin=2.5cm}

\begin{document}

\section*{Domain Agnostic Fourier Neural Operators}
\textbf{Autoren:} Ning Liu, Siavash Jafarzadeh, Yue Yu \\
\textbf{Konferenz:} NeurIPS 2023 \\
\textbf{Code:} \url{https://github.com/ningliu-iga/DAFNO}

% -------------------------------------------------------------------
\subsection*{Problemstellung und Motivation}

Standard-Fourier Neural Operators (FNOs) sind leistungsf\"ahige Surrogat\-modelle
f\"ur partielle Differentialgleichungen, die durch Lernen im Spektralbereich mittels
der Fast Fourier Transform (FFT) hohe Recheneffizienz erreichen. Ihre zentrale
Schwachstelle liegt jedoch in der Anforderung an die Eingabegeometrie: Die FFT
setzt ein rechteckiges Rechengebiet mit gleichm\"a\ss{}igem Gitter voraus. In der
Praxis sind physikalische Gebiete aber h\"aufig unregelm\"a\ss{}ig geformt -- zum
Beispiel Materialproben mit Hohlr\"aumen, Fl\"ugelprofilen oder sich ausbreitenden
Rissen. Der ga\"ngige Ausweg, das Gebiet durch Nullauff\"ullung auf ein Rechteck zu
erweitern, f\"uhrt zu numerischen Instabilit\"aten (Gibbs-Ph\"anomen) und versagt
v\"ollig bei Gebieten mit topologischen \"Anderungen, etwa wachsenden Rissen.

% -------------------------------------------------------------------
\subsection*{Kernidee: Domain Agnostic FNO (DAFNO)}

Die Autoren schlagen vor, die Gebietsinformation explizit in die Architektur der
FNO-Schichten einzubetten, anstatt sie als nachgelagerte Korrektur zu behandeln.
Konkret wird eine \emph{charakteristische Funktion} des Rechengebiets~$\Omega$
definiert:
\[
    \chi(\mathbf{x}) =
    \begin{cases}
        1 & \mathbf{x} \in \Omega \\
        0 & \mathbf{x} \in \mathbb{T} \setminus \Omega
    \end{cases}
\]
Dabei ist $\mathbb{T}$ eine das Gebiet~$\Omega$ umschlie\ss{}ende rechteckige
Box. Diese Funktion wird in den Faltungskern der FNO-Schicht eingebaut, so dass
die Aktualisierungsregel eines DAFNO-Layers lautet:
\[
    \mathcal{J}^l[\mathbf{h}]
    = \sigma\!\left(
        \chi(\mathbf{x})\!\left[
            \mathcal{I}\!\left(\chi(\cdot)\mathbf{h}(\cdot);\mathbf{v}^l\right)
            - \mathbf{h}(\mathbf{x})\,\mathcal{I}\!\left(\chi(\cdot);\mathbf{v}^l\right)
            + W^l \mathbf{h}(\mathbf{x}) + \mathbf{c}^l
        \right]
    \right),
\]
wobei $\mathcal{I}(\circ;\mathbf{v}^l) := \mathcal{F}^{-1}[\mathcal{F}[\kappa(\cdot;\mathbf{v}^l)]
\cdot \mathcal{F}[\circ]]$ den \"ublichen Faltungsoperator mit FFT bezeichnet.
Da der Integrationsbereich weiterhin die rechteckige Box~$\mathbb{T}$ umfasst,
bleibt die FFT anwendbar und die Effizienz erhalten. Das Produkt
$\chi(\mathbf{x})\chi(\mathbf{y})$ sorgt daf\"ur, dass Interaktionen zwischen
Punkten innerhalb und au\ss{}erhalb des physikalischen Gebiets unterdr\"uckt werden.

Da die scharfe charakteristische Funktion~$\chi$ nicht stetig ist und damit das
Gibbs-Ph\"anomen hervorruft, wird sie durch eine gegl\"attete Variante ersetzt:
\[
    \tilde{\chi}(\mathbf{x})
    := \tanh\!\left(\beta \cdot \mathrm{dist}(\mathbf{x}, \partial\Omega)\right)
       \left(\chi(\mathbf{x}) - 0{,}5\right) + 0{,}5.
\]
Der Parameter~$\beta$ steuert den \"Ubergang zwischen Gebiet und Rand. Diese
Gl\"attung verbessert die spektrale Repr\"asentation erheblich und reduziert
numerische Artefakte nahe der Gebietsgrenze deutlich.

% -------------------------------------------------------------------
\subsection*{Architekturvarianten}

Die Autoren stellen zwei Varianten vor:
\begin{itemize}
    \item \textbf{eDAFNO} (explicit): basierend auf dem Standard-FNO mit
          schichtabh\"angigen Parametern; leicht genauer, h\"oherer
          Speicherbedarf.
    \item \textbf{iDAFNO} (implicit): basierend auf dem IFNO mit
          schichtunabh\"angigen Parametern; deutlich kompakter (nur ein Viertel
          der Parameter), weniger anf\"allig f\"ur \"Uberanpassung bei kleinen
          Datens\"atzen.
\end{itemize}

% -------------------------------------------------------------------
\subsection*{Experimente und Ergebnisse}

Die Methode wird an drei unterschiedlichen Benchmarks validiert:

\paragraph{Hyperelastizit\"at.}
Vorhersage von Spannungsfeldern in Materialproben mit zuf\"allig geformten
Hohlr\"aumen. DAFNO erzielt bei 1000~Trainingssamples einen relativen
$L_2$-Fehler von 1,094\,\%, w\"ahrend das beste nicht-DAFNO-Baseline-Modell
(Geo-FNO) 2,316\,\% erreicht -- eine Verbesserung um \"uber 100\,\%.

\paragraph{Fl\"ugelprofil-Design.}
Vorhersage des Mach-Zahlen-Felds um NACA-0012-Varianten. eDAFNO erreicht
0,596\,\% Testfehler und \"ubertrifft Geo-FNO (1,650\,\%) um 63,9\,\%.

\paragraph{Rissausbreitung.}
In diesem Experiment wird die besondere St\"arke des DAFNO deutlich:
Mit nur einem einzigen Trainingsszenario ($\sigma = 4\,\mathrm{MPa}$) kann
das Modell Rissverl\"aufe unter v\"ollig verschiedenen Lasten ($\sigma = 2$ und
$6\,\mathrm{MPa}$) -- inklusive Rissver\"astelung und Mehrfachrissen -- korrekt
vorhersagen. Standard-FNO mit Maske scheitert hier bereits nach kurzer Simulationszeit,
da die entstehenden Geometrien au\ss{}erhalb der Trainingsverteilung liegen.

% -------------------------------------------------------------------
\subsection*{Grenzen}

In der aktuellen Implementierung setzt DAFNO gleichm\"a\ss{}ige Gitter voraus.
F\"ur unregelm\"a\ss{}ige Gitter ben\"otigt man eine zus\"atzliche analytische
oder trainierbare Koordinatentransformation, wie im Fl\"ugelprofil-Experiment
gezeigt. Dar\"uber hinaus ist der Transfer zwischen verschiedenen Typen von
Randbedingungen (z.\,B. Dirichlet vs.\ Neumann) noch nicht untersucht.

% -------------------------------------------------------------------
\subsection*{Bezug zur Masterarbeit}

Die Kernintuition von DAFNO -- explizite geometrische Information durch eine
masken\"ahnliche Funktion in die Netzwerkarchitektur einzubetten -- findet sich
in der vorliegenden Masterarbeit in analoger Form wieder. Der in dieser Arbeit
verwendete vierkan\"alige Eingabetensor enth\"alt als ersten Kanal eine bin\"are
Kristallmaske $M(\mathbf{x},z) \in \{0,1\}$, die genauso wie~$\chi$ die
Positionen der festen Einschl\"usse (Kristalle) auf dem regelm\"a\ss{}igen
$256\times256$-Gitter markiert. Erg\"anzt wird diese durch ein vorzeichenbehaftetes
Distanzfeld als zweiten Kanal, das einer gegl\"atteten Version von~$\chi$
konzeptionell verwandt ist: beide Repr\"asentationen \"uberf\"uhren die scharfe
bin\"are Geometriebeschreibung in eine stetigere, f\"ur die Spektralanalyse
g\"unstigere Form.

Der zentrale Unterschied liegt in der \emph{Tiefe der Einbettung}: W\"ahrend
DAFNO die Gebietsinformation durch Multiplikation mit~$\chi$ direkt in jeden
Fourier-Layer integriert und damit die Interaktion zwischen Punkten innerhalb
und au\ss{}erhalb des Gebiets auf Architekturebene kontrolliert, flie\ss{}t die
Maske in der vorliegenden Arbeit als zus\"atzlicher Eingabekanal in das Netzwerk
ein und wird durch das Training implizit ber\"ucksichtigt. Die DAFNO-Methode
stellt damit eine architekturelle Versch\"arfung des in dieser Arbeit genutzten
Ansatzes dar, die potenziell pr\"azisere Randbedingungen an den Kristalloberfl\"achen
erm\"oglichen k\"onnte.

F\"ur das spezifische Problem der vorliegenden Arbeit -- station\"are Stokes-Str\"omung
um starre, kreisf\"ormige Kristalle auf einem festen Gitter -- ist der einfachere
Eingabekanal-Ansatz gut motiviert: Da sich die Kristallgeometrie selbst nicht \"andert
und die Netzwerke ausschlie\ss{}lich auf \"aquidistanten Gittern trainiert und ausgewertet
werden, entfallen die topologischen \"Anderungen und die Probleme mit Gitterirregularit\"aten,
f\"ur die DAFNO prim\"ar konzipiert wurde. Die Ergebnisse des DAFNO-Papers best\"atigen
dennoch die grunds\"atzliche Herangehensweise: Die explizite geometrische Kodierung
der Einschl\"usse ist ein entscheidender Baustein f\"ur gute Generalisierung \"uber
verschiedene r\"aumliche Konfigurationen hinweg.

Dar\"uber hinaus ist die in DAFNO gezeigte F\"ahigkeit zur Generalisierung \"uber
topologisch verschiedene Gebiete -- mit nur wenigen Trainingssimulationen -- direkt
relevant f\"ur die Forschungsfrage der Masterarbeit: Kann ein auf $N$ Kristallen
trainiertes Modell auf Konfigurationen mit $1$ bis $N+m$ Kristallen generalisieren?
Die DAFNO-Ergebnisse legen nahe, dass eine explizite geometrische Kodierung die
Out-of-distribution-Generalisierung erheblich verbessert -- eine Beobachtung, die
die Nutzung der Kristallmaske als zentralen Eingabekanal in der vorliegenden
Arbeit nachtr\"aglich st\"utzt und als Motivation f\"ur zuk\"unftige Erweiterungen
in Richtung DAFNO-artiger Integrationsstrategien dienen kann.

\end{document}